% Block diagram of the firmware algorithm on the ESP32-S3.
\begin{tikzpicture}[
    font=\footnotesize,
    node distance=0.36cm,
    term/.style={draw, thick, rounded rectangle, align=center, minimum height=0.8cm, text width=4.2cm, fill=gray!10},
    proc/.style={draw, thick, align=center, minimum height=0.8cm, text width=6.2cm, fill=white},
    io/.style={draw, thick, trapezium, trapezium left angle=72, trapezium right angle=108, trapezium stretches=true,
               align=center, minimum height=0.8cm, text width=5.2cm, fill=blue!5},
    dec/.style={draw, thick, diamond, aspect=2.4, align=center, inner sep=1pt, text width=2.6cm, fill=yellow!10},
    arr/.style={-{Stealth[length=2.2mm]}, thick},
  ]
  \node[term] (start) {Старт / захранване};
  \node[proc, below=of start] (init) {Инициализация: тактова честота,\\SPI, GPIO, радио (BLE/Wi-Fi)};
  \node[proc, below=of init] (cfg) {Конфигурация на ADS1292R: 500 SPS,\\PGA = 6, RLD вкл., канал 2 активен};
  \node[proc, below=of cfg] (load) {Зареждане на тегла, \(s\), SOS и int8\\параметри; нулиране на филтрите};
  \node[proc, below=of load] (go) {START = 1, разрешаване на\\прекъсване по DRDY};
  \node[io, below=of go] (rd) {DRDY \(\downarrow\): четене на 9 байта\\(статус + 2 \(\times\) 24 bit)};
  \node[dec, below=of rd] (lo) {Lead-off?};
  \node[proc, right=0.9cm of lo, text width=3.0cm] (flag) {Флаг \enquote{лош контакт} в следващия пакет};
  \node[proc, below=of lo] (mv) {Код \(\rightarrow\) mV;\; НЧ FIR филтър\\и децимация \(5{:}1\) (500 \(\rightarrow\) 100 Hz)};
  \node[dec, below=of mv] (dec5) {Нов отчет\\при 100 Hz?};
  \node[proc, below=of dec5] (hp) {Каузален ВЧ филтър (SOS);\\запис в кръгов буфер};
  \node[dec, below=of hp] (full) {256 нови\\отчета?};
  \node[proc, below=of full] (enc) {\(x/s\) \(\rightarrow\) енкодер \(\rightarrow z\) (float32)\\\(\rightarrow\) int8 квантизация};
  \node[io, below=of enc] (tx) {Пакет: №, флагове, 32 B \(z\)\\\(\rightarrow\) изпращане};

  \draw[arr] (start) -- (init); \draw[arr] (init) -- (cfg); \draw[arr] (cfg) -- (load);
  \draw[arr] (load) -- (go); \draw[arr] (go) -- (rd); \draw[arr] (rd) -- (lo);
  \draw[arr] (lo) -- node[above] {да} (flag);
  \draw[arr] (lo) -- node[right] {не} (mv);
  \draw[arr] (flag) |- (mv);
  \draw[arr] (mv) -- (dec5);
  \draw[arr] (dec5) -- node[right] {да} (hp);
  \draw[arr] (hp) -- (full);
  \draw[arr] (full) -- node[right] {да} (enc);
  \draw[arr] (enc) -- (tx);
  % loops back to waiting for the next DRDY
  \coordinate (back) at ($(rd.west)+(-1.6,0)$);
  \draw[arr] (dec5.west) -- node[above, pos=0.3] {не} (dec5.west -| back) -- (back) -- (rd.west);
  \draw[arr] (full.west) -- node[above, pos=0.3] {не} (full.west -| back) -- (back);
  \draw[arr] (tx.west) -- (tx.west -| back) -- (back);
\end{tikzpicture}
